\section{Introduction}

Multimodal Large Language Models (MLLMs) and diffusion-based generative models~\cite{} have each achieved remarkable progress in recent years. Despite their respective successes, these two model families have evolved largely in isolation: one focuses on interpreting the visual world, and the other on generating it. This naturally motivates the constrcution of a unified model capable of jointly performing understanding and generation. In response, recent Recent Unified Multimodal Models (UMMs)~\cite{} have emerged as a early attempts to integrate these capabilities within a single architecture, achieving performance comparable to expert models in each domain. However, this progress leads to a fundamental inquiry: is the goal of UMMs simply to consolidate two expert systems into one parameter space, or to achieve genuine synergy between understanding and generation that enables a deeper form of world modeling?

% To address this question, we must first determine whether meaningful synergy actually emerges within UMMs. 
Several prior works suggest that complementarity between perception and generation is possible. For example, REPA improves diffusion generation by aligning intermediate DiT representations with DINOv2 features, and other works have shown that auxiliary reconstruction objectives can enhance visual understanding in MLLMs. However, it remains unclear whether such synergy naturally arises within unified architectures. While some studies report positive interactions between tasks, others observe interference, where joint modeling degrades performance. We argue that this discrepancy stems from how synergy has been evaluated. Existing approaches typically rely on aggregate performance changes on general multimodal benchmarks, which do not directly measure whether knowledge learned in one task influences the representations used for another. True synergy cannot be reduced to parameter sharing; it requires knowledge integration, where knowledge acquired from one task meaningfully shapes the representations underlying the other. Therefore, evaluating synergy in UMMs demands a more direct measurement of cross-task knowledge transfer.

To this end, we propose task transferability as a concrete indicator of knowledge integration. Specifically, we examine whether learning in one modality, understanding or generation, can directly improve performance in the other without explicit supervision in that target modality.  As a testbed, we focus on counting, a capability known to be challenging for both MLLMs and diffusion models. We measure whether counting-focused visual instruction tuning improves counting-aware image generation, and conversely, whether counting-aware generative training enhances counting performance in understanding. Through systematic experiments across representative UMM architectures, we find a clear architectural distinction between models that exhibit bi-directional transferability and those that do not. Transferability consistently emerges only in architectures equipped with a shared transformer backbone and a unified image tokenizer. In contrast, loosely coupled designs fail to demonstrate meaningful cross-modal improvement. These findings suggest that knowledge integration is not an automatic consequence of joint training, but is fundamentally shaped by architectural design, highlighting the importance of structural coupling for enabling synergy.

Finally, we extend our analysis beyond counting to investigate how transferability manifests across tasks of varying granularity. From fine-grained tasks such as OCR to coarse semantic reasoning such as spatial position estimation, we observe that different tasks exhibit modality preference—some are more effectively learned through understanding and then transferred to generation, while others benefit from the reverse direction. Notably, in several cases, training in the modality more naturally aligned with a task and subsequently transferring to the target modality yields greater efficiency than direct supervision in the target itself. These findings reveal that the relationship between understanding and generation is neither trivial nor symmetric, but depends on both task characteristics and architectural structure. Our study thus reframes UMM design not as simple functional unification, but as the deliberate construction of architectures capable of genuine knowledge integration

Our contributions ar as follows:
\begin{itemize}
    \item first
    \item second
    \item third
\end{itemize}